% \iffalse meta-comment
%
% Copyright (C) 2017- by Yanshuo Chu <yanshuoc@gmail.com>
%
% This file may be distributed and/or modified under the
% conditions of the LaTeX Project Public License, either version 1.3a
% of this license or (at your option) any later version.
% The latest version of this license is in:
%
% http://www.latex-project.org/lppl.txt
%
% and version 1.3a or later is part of all distributions of LaTeX
% version 2004/10/01 or later.
%
% \fi
%
% \section{终稿的脚注}
% \changes{v3.2a}{2026/09/26}{终稿的脚注交给 \pkg{docxlike} 照 Word 排，不再加载 \pkg{footmisc}：
%   脚注号上标带圈、每页重编，脚注文字小五单倍，分隔线 144 磅，去掉没有出处的 1.5 字悬挂}
% \changes{v3.2a}{2026/09/13}{脚注整层挪到 \file{src/engine/hit-footnote.dtx} 共用，这里只剩
%   套 \option{stage} 闸的调用}
%
% ^^A ======================================================================
% ^^A Module final-footnote: 脚注样式与编号（hit-thesis）
% ^^A ======================================================================
%    \begin{macrocode}
%<*final-footnote>
\bool_if:NT \g__hit_final_bool
  {
%    \end{macrocode}
%
% 脚注样式见 \file{src/engine/hit-footnote.dtx}，那里定义的带圈数字命令（\cs{quan} 等）与
% \cs{footnotesize} 照用。脚注本身交给 \pkg{docxlike} 照 Word 排：编号带圈、每页重新编号（人文社科类
% 写作指南“页下以①②③……来表示”，理工类没写，一并照这样），脚注号是 \texttt{footnote reference}
% 上标，脚注文字是 \texttt{footnote text}（小五、单倍、不贴行网格，Word 的内建定义），分隔线是
% Normal 段基线上 144 磅长的线，脚注排在页底。先前脚注号后悬挂 1.5 个字、分隔线 0.3 版心宽，
% 现行写作指南里都找不到出处，改照 Word 插脚注的样子。编号的字仍交给 \cs{quan}：宋体只有 ①--⑩，
% 再往后字体里没有，\cs{quan} 会画圈补上（Word 是换字体补字）。先前的版本借 \pkg{footmisc} 的
% \texttt{perpage}，\file{aux} 里留着 \pkg{perpage} 的 \cs{pp@pagectr}，这里给个空定义，升级后
% 第一遍读旧 \file{aux} 不报错。
%    \begin{macrocode}
\hit@setup@footnote
\docxlikesetup
  { footnotes = { number-format = decimal-enclosed-circle , numbering = restart-each-page } }
\cs_gset:Npn \thefootnote { \quan { \docxlike_footnote_index: } }
\cs_if_exist:NF \pp@pagectr { \cs_new:Npn \pp@pagectr #1#2#3#4 { } }
  }
%</final-footnote>
%    \end{macrocode}
%
% \Finale
